\documentclass[UTF8,11pt]{ctexart}
\usepackage[a4paper,margin=25mm]{geometry}
\usepackage{amsmath,amssymb,amsthm,mathtools}
\usepackage[all]{xy}
\usepackage{xurl}
\usepackage{hyperref}
\usepackage{enumitem}
\hypersetup{hidelinks,breaklinks=true}
\setlength{\emergencystretch}{3em}
\newtheorem*{theorem}{Theorem}
\newtheorem*{lemma}{Lemma}
\newtheorem*{proposition}{Proposition}
\newtheorem*{corollary}{Corollary}
\newtheorem*{definition}{Definition}
\newcommand{\Spec}{\operatorname{Spec}}
\newcommand{\Hom}{\operatorname{Hom}}
\newcommand{\End}{\operatorname{End}}
\newcommand{\Gal}{\operatorname{Gal}}
\newcommand{\Cl}{\operatorname{Cl}}
\newcommand{\vol}{\operatorname{vol}}
\newcommand{\covol}{\operatorname{covol}}
\newcommand{\tr}{\operatorname{tr}}
\newcommand{\im}{\operatorname{im}}
\newcommand{\id}{\operatorname{id}}
\newcommand{\coker}{\operatorname{coker}}
\newcommand{\stacksref}[2]{\href{https://stacks.math.columbia.edu/tag/#1}{#2 [#1]}}
\begin{document}
\section*{008\quad 三维直线的joints数量界}
\subsection*{来源与整理范围}
集体来源：MIT OpenCourseWare，MIT 18.S997 \emph{The Polynomial Method}，Fall 2012；授课教师 Lawrence D. Guth。课程页面注明讲义由 MIT 学生提供并获准使用，所取文件未单列执笔人，故不将转录正文擅自署为 Guth 独著。本文收录的是该课程的人类讲义版本，不是研究论文原版。获取日期：2026-09-28。\par 原文：Lecture 4, \emph{The Joints Problem}, Theorem 0.1、Main Lemma、Lemma 0.2 及所含全部证明，pp.~1--2。维数计数引用 Lecture 1, Corollaries 0.2--0.3，pp.~1--2。
\par\url{https://ocw.mit.edu/courses/18-s997-the-polynomial-method-fall-2012/a6ce9ea12a24f8e8156cbf0857f7e13e_MIT18_S997F12_lec4.pdf}
\par\url{https://ocw.mit.edu/courses/18-s997-the-polynomial-method-fall-2012/43c5d5700347be070289fe2be4cf30c7_MIT18_S997F12_lec1.pdf}
\par\url{https://ocw.mit.edu/courses/18-s997-the-polynomial-method-fall-2012/pages/lecture-notes/}

\subsection*{作者原命题与人类证明}
Suppose that we have a set of lines in $\mathbb R^3$. A joint of the set of lines is a point which lies in three non-coplanar lines. The joints problem asks what is the maximal number of joints that can be formed from $L$ lines.

\begin{theorem}[0.1]
Any $L$ lines in space determine $\leq 10L^{3/2}$ joints.
\end{theorem}
\begin{lemma}[Main Lemma]
If a set of lines has $J$ joints, then one of the lines contains $\leq 3J^{1/3}$ joints.
\end{lemma}
The main lemma implies the theorem by removing the lines one at a time. We start with $L$ lines and $J$ joints. By cutting out one line, we reduce the number of joints by $\leq 3J^{1/3}$. We look at the remaining lines $L-1$ lines, which contain $\leq J$ joints. One of the lines has $\leq 3J^{1/3}$ joints on it. Removing this line, we reduce the number of joints by $\leq 3J^{1/3}$. We remove all $L$ lines, one at a time. Each time the number of joints decreases by $\leq 3J^{1/3}$, and we end up with no joints. Therefore, $J\leq L(3J^{1/3})$. Rearranging we get $J^{2/3}\leq 3L$, which implies the theorem.

Now we turn to the proof of the main lemma.
\begin{proof}
Let $P$ be a lowest degree non-zero polynomial that vanishes at every joint. The degree of $P$ is $\leq 3J^{1/3}$ by dimension counting, as in Lecture 1. If every line has $>3J^{1/3}$ joints, then $P$ must vanish on every line.

\begin{lemma}[0.2]
If $x$ is a joint lying in three (non-coplanar) lines, and if a smooth function $F:\mathbb R^3\to\mathbb R$ vanishes on the lines, then $\nabla F$ vanishes at $x$.
\end{lemma}
\begin{proof}
Let $v_1,v_2,v_3$ be tangent vectors for the three lines. The directional derivative of $F$ in the direction $v_i$ must vanish at $x$. So we have $\nabla F(x)\cdot v_i=0$ for each $i$. Since the $v_i$ are a basis of $\mathbb R^3$, we have $\nabla F(x)=0$.
\end{proof}
So we see that the derivates of $P$ vanish at each joint. The derivates have smaller degree than $P$. Since $P$ was a minimal degree non-zero polynomial that vanishes at each joint, the derivatives of $P$ are all the zero polynomial! Then $P$ must be constant, and we get a contradiction.
\end{proof}


\subsection*{整理说明（非作者正文）}
保留作者先用主引理推出定理、再证明主引理的顺序，嵌入的 Lemma 0.2 及其证明未略去。只略去原文的历史介绍、例子与主证明后的网格例子；这些不是该定理证明的续证。

标准先修：一元多项式的根数界；Lecture 1 的维数计数（原文如下）。$V(d)$ 为 $n$ 元总次数至多 $d$ 的多项式空间，$s$ 为点数。没有把 joints 定理当作先修。

原文重复的 ``remaining lines $L-1$ lines'' 和拼写 ``derivates'' 保留。取得了 PDF 文本，但截图接口失败，未声称逐页图像核对或保留了下载原件。数学公式按实际取得文本转录。

\paragraph{所引用的维数计数结果（Lecture 1，pp.~1--2）}
\begin{corollary}[0.2]
If $\dim V(d)>s$, then there is a non-zero polynomial $P$ of degree $\leq d$ that vanishes on the given finite set.
\end{corollary}
The dimension of $V(d)$ is $\binom{d+n}{n}\geq d^n/n!$. For $n$ fixed and $d$ large, $d^n/n!$ is a good approximation of the dimension. Therefore, we get the following corollary.
\begin{corollary}[0.3]
For any set of $s$ points in $F^n$, there is a non-zero polynomial that vanishes on the set with degree $\leq ns^{1/n}$.
\end{corollary}


\subsection*{可修改的简短标注（非作者正文）}
$c$：三维直线配置的交会点极值。

$a$：最小次数消失多项式；维数计数；直线限制与根数界；方向导数与梯度；最小次数反证；逐线删除计数。

\texttt{cross\_theory\_status = yes}。三条不共面直线的几何条件转成梯度对三个独立方向同时正交；偏导数降次与最小次数性产生一条稀疏直线，再返回组合删除计数。位置：Main Lemma 与 Lemma 0.2，p.~2。全部标注为非穷尽、可修改初稿。
\subsection*{许可与修改记录}
材料来自 MIT OpenCourseWare；原作者与课程见上文。按该站所示 \href{https://creativecommons.org/licenses/by-nc-sa/4.0/}{CC BY-NC-SA 4.0} 许可复用。本转录及其附加整理采用同一许可；不得据此声称作者或 MIT 认可本次整理。2026-09-28：ChatGPT 转录为 TeX，加入题号、来源、整理说明和初稿标注；数学正文保持作者语言及顺序。

\end{document}
